\documentclass[11pt]{article}

\usepackage[margin=1in]{geometry}
\usepackage{amsmath}
\usepackage{amssymb}
\usepackage{booktabs}
\usepackage[T1]{fontenc}
\usepackage[utf8]{inputenc}
\usepackage{lmodern}

\title{How \texttt{reactorModel.ipynb} Drives the Reactor Web App}
\author{}
\date{}

\begin{document}

\maketitle

\section{Purpose}

This note explains how the repository's physics is split across:
\begin{itemize}
  \item the exploratory notebook \texttt{theory/reactorModel.ipynb},
  \item the point-kinetics web simulation shown on the \emph{Overview} page,
  \item the steady 2D diffusion solve shown on the \emph{Core} page,
  \item and the direct time-dependent diffusion solve shown on the
  \emph{Transient Diffusion} page.
\end{itemize}

The notebook is the \textbf{source of modeling logic and calibration workflow}.
The web app is the \textbf{interactive execution layer}.  The most important
architectural point is:

\begin{quote}
The repository now contains \textbf{two transient models}: a fast
\textbf{point-kinetics surrogate} for the Overview page, and a
\textbf{spatially resolved 2D diffusion transient} for the Transient Diffusion
page.  Both are tied back to the same one-group annular-core assumptions and
the same notebook-derived calibration data.
\end{quote}

\section{What the notebook does}

The notebook \texttt{theory/reactorModel.ipynb} studies a heavy-water moderated
annular reactor with cylindrical symmetry in $(r,z)$ and a top-inserted central
control/shutdown bank.  It has two estimation levels.

\subsection{Estimate 1: homogeneous buckling sanity check}

The first estimate is a coarse analytical model based on one-speed diffusion
theory in a homogeneous bare cylinder:
\[
-D\nabla^2 \phi + \Sigma_a \phi = \frac{1}{k}\nu\Sigma_f \phi.
\]

The criticality condition is written as
\[
B_m^2 = \frac{\nu\Sigma_f - \Sigma_a}{D} = B_g^2,
\]
with geometric buckling for a finite cylinder approximated by
\[
B_g^2 \approx \left(\frac{2.405}{R}\right)^2 + \left(\frac{\pi}{H}\right)^2.
\]

In the notebook this estimate is used only as a \textbf{first-pass size check}.
It is intentionally approximate because it ignores:
\begin{itemize}
  \item the inner moderator hole,
  \item reflector heterogeneity,
  \item explicit axial reflector slabs,
  \item and the geometric control rod.
\end{itemize}

\subsection{Estimate 2: full 2D $r$--$z$ finite-difference diffusion model}

The second estimate is the notebook's main result.  It solves a one-group
diffusion eigenvalue problem on a cylindrical $r$--$z$ mesh:
\[
-\nabla \cdot \left(D \nabla \phi\right) + \Sigma_a \phi
= \frac{1}{k_{\mathrm{eff}}}\nu\Sigma_f \phi.
\]

This estimate captures:
\begin{itemize}
  \item radial and axial leakage,
  \item the annular fuel geometry,
  \item inner and outer D$_2$O regions,
  \item axial reflector slabs,
  \item and the real rod tip position as insertion changes.
\end{itemize}

The notebook uses this model to:
\begin{enumerate}
  \item search for critical annular geometry,
  \item compute the clean-core eigenvalue,
  \item scan rod insertion versus reactivity,
  \item and convert that scan into the calibration constants later frozen in the
        backend.
\end{enumerate}

\section{Geometry, materials, and assumptions}

\subsection{Core layout}

The modeled reactor is an annular cylindrical core with a central D$_2$O region,
a homogenized fuel annulus, and an outer D$_2$O reflector.  The current frozen
geometry in \texttt{backend/reactor\_backend/}\linebreak\texttt{calibration.py}
is:

\begin{center}
\begin{tabular}{lc}
\toprule
Quantity & Value \\
\midrule
Inner moderator radius $R_{\mathrm{inner}}$ & $80.0$ cm \\
Fuel outer radius $R_{\mathrm{fuel}}$ & $345.6$ cm \\
Reflector outer radius $R_{\mathrm{refl}}$ & $405.6$ cm \\
Active height $H$ & $691.2$ cm \\
Axial reflector thickness $H_{\mathrm{refl}}$ & $60.0$ cm per side \\
Equivalent rod radius $r_{\mathrm{rod}}$ & $50.0$ cm \\
\bottomrule
\end{tabular}
\end{center}

The model is axisymmetric, so the mesh covers only $r \ge 0$ while still
representing the full 3D core through cylindrical volume weighting
$2\pi r\,dr\,dz$.

\subsection{One-group effective constants}

The notebook starts from microscopic cross sections at the 2200 m/s thermal
reference point and then builds \textbf{engineering-effective} one-group
constants.  The production code uses the following frozen values:

\begin{center}
\begin{tabular}{lccc}
\toprule
Region & $D$ [cm] & $\Sigma_a$ [cm$^{-1}$] & $\nu\Sigma_f$ [cm$^{-1}$] \\
\midrule
Inner D$_2$O moderator & $1.25$ & $3.5\times 10^{-4}$ & $0$ \\
Fuel annulus & $0.95$ & $8.5\times 10^{-3}$ & $8.59\times 10^{-3}$ \\
Outer D$_2$O reflector & $1.15$ & $4.0\times 10^{-4}$ & $0$ \\
\bottomrule
\end{tabular}
\end{center}

These are \textbf{homogenized one-group constants}, not a multigroup transport
library.  The fuel values intentionally fold together:
\begin{itemize}
  \item natural uranium fuel,
  \item D$_2$O moderation inside the lattice,
  \item resonance/self-shielding effects,
  \item and spectrum averaging into one effective thermal group.
\end{itemize}

The control bank is represented as an added absorber in the inner moderator
region with
\[
\Delta \Sigma_{a,\mathrm{rod,max}} = 0.25\ \mathrm{cm}^{-1}.
\]
This is likewise an effective one-group quantity rather than a detailed rod-cell
transport model.

\subsection{How the physics was manipulated into these constants}

This project deliberately compresses more detailed reactor physics into a
one-group annular model that can run interactively.  The main reductions are:
\begin{enumerate}
  \item \textbf{Energy collapse.}  Multigroup slowing-down and spectral effects
        are represented by a single thermal group with effective
        $(D,\Sigma_a,\nu\Sigma_f)$.
  \item \textbf{Lattice homogenization.}  The fuel annulus is not modeled as
        explicit fuel pins in moderator.  Instead natural uranium and D$_2$O are
        volume-averaged into one fissile ring region.
  \item \textbf{Leakage folded into effective losses.}  In a one-group model,
        any loss of neutrons to higher-energy groups or to unresolved geometric
        detail is represented as part of an effective absorption-like loss.
  \item \textbf{Equivalent rod model.}  The real control/shutdown system is
        collapsed to one central cylindrical absorber with an effective radius
        and an effective extra absorption $\Delta\Sigma_a$.
  \item \textbf{Calibration by 2D diffusion, not by transport.}  The final rod
        worth curve and critical insertion are taken from repeated diffusion
        eigenvalue solves, then frozen for the point model.
\end{enumerate}

The notebook is explicit about these manipulations.  For example, the D$_2$O
``effective'' $\Sigma_a$ is much larger than the true 2200 m/s absorption of
pure heavy water, because the one-group model uses $\Sigma_a$ as a \emph{net
loss coefficient} after spectral and geometric simplification.  Likewise, the
fuel-annulus $\Sigma_a$ is intentionally inflated relative to a naive
2200 m/s natural-uranium homogenization so the one-group annular model better
matches the critical-size and rod-worth behavior seen in the notebook's 2D
studies.

\subsection{Boundary assumptions}

The diffusion solves use extrapolated zero-flux boundaries rather than placing
the computational boundary exactly at the physical reflector edge.  In code this
is represented by an extrapolation factor of $2.13 D_{\mathrm{refl}}$, so the
radial and axial computational extents are larger than the physical core.

\section{How the notebook feeds the web app}

The notebook does \emph{not} run the web app directly.  Instead it establishes
the physical basis that is later frozen into backend constants and reusable
solver code:

\begin{enumerate}
  \item \textbf{Estimate 1} gives an analytical reference for rough critical size.
  \item \textbf{Estimate 2} performs the annular $r$--$z$ diffusion study that
        supplies the production calibration.
  \item The resulting geometry, material constants, and rod-worth data are copied
        into \texttt{calibration.py}.
  \item The backend then exposes two kinds of runtime models:
        \begin{itemize}
          \item a calibrated point-kinetics model for fast dashboard transients,
          \item and spatial diffusion solvers for the Core and Transient pages.
        \end{itemize}
\end{enumerate}

So the notebook is best viewed as the \textbf{physics derivation and calibration
workspace}, while the Python backend is the \textbf{production implementation}.

\subsection{How the notebook output becomes backend input}

The handoff from notebook to code is not symbolic; it is numerical.  The
production backend lifts the notebook's final values into
\texttt{calibration.py}:
\begin{itemize}
  \item annular geometry dimensions,
  \item one-group material constants,
  \item clean-core $k_{\mathrm{eff}}$ and excess reactivity,
  \item the 11-point rod-worth table,
  \item and the derived critical insertion percentage.
\end{itemize}

This means the backend no longer recomputes those design-calibration studies on
startup.  Instead it assumes the notebook has already answered the
``what reactor are we modeling?'' question and the Python solvers answer the
``how does that reactor evolve for this input?'' question.

\section{Overview page: point-kinetics transient model}

\subsection{Governing equations}

The Overview page uses the standard six-group delayed-neutron point-kinetics
system:
\begin{align}
\frac{dn}{dt} &= \frac{\rho - \beta}{\Lambda} n + \sum_{i=1}^{6}\lambda_i C_i,
\\
\frac{dC_i}{dt} &= \frac{\beta_i}{\Lambda} n - \lambda_i C_i,
\qquad i=1,\dots,6.
\end{align}

Here:
\begin{itemize}
  \item $n(t)$ is the normalized neutron population,
  \item $\rho$ is reactivity in $\Delta k/k$,
  \item $\beta = \sum_i \beta_i$ is the effective delayed-neutron fraction,
  \item $\Lambda$ is the prompt neutron generation time,
  \item $C_i$ is precursor concentration for delayed group $i$,
  \item $\lambda_i$ is the delayed-group decay constant.
\end{itemize}

\subsection{Constants used by the web app}

The current production constants in \texttt{config.py} are:
\begin{align*}
\beta_{\mathrm{eff}} &= 0.00651 = 651\ \mathrm{pcm}, \\
\Lambda &= 5\times 10^{-4}\ \mathrm{s}.
\end{align*}

The six delayed groups are:

\begin{center}
\begin{tabular}{ccc}
\toprule
Group & $\beta_i$ & $\lambda_i$ [s$^{-1}$] \\
\midrule
1 & 0.00025 & 0.0124 \\
2 & 0.00138 & 0.0305 \\
3 & 0.00122 & 0.1110 \\
4 & 0.00264 & 0.3010 \\
5 & 0.00075 & 1.1400 \\
6 & 0.00027 & 3.0100 \\
\bottomrule
\end{tabular}
\end{center}

Displayed output is scaled from the nominal operating point
\[
P_0 = 20\ \mathrm{MW_{th}}, \qquad
\Phi_0 = 1.5\times 10^{12}\ \mathrm{n\,cm^{-2}\,s^{-1}},
\]
using
\[
P(t) = P_0 n(t), \qquad \Phi(t) = \Phi_0 n(t).
\]

\subsection{How rod motion enters the point model}

The point-kinetics solver does not compute rod worth from first principles during
runtime.  Instead it reuses the \textbf{2D diffusion-derived rod-worth table}
frozen in \texttt{calibration.py}.  Reactivity is computed as
\[
\rho_{\mathrm{tot,pcm}}(x)
= \rho_{\mathrm{clean,pcm}} + \Delta \rho_{\mathrm{rod,pcm}}(x)
+ \rho_{\mathrm{scram,pcm}},
\]
where $x \in [0,1]$ is rod insertion fraction.  In the current calibration:
\begin{align*}
k_{\mathrm{clean}} &= 1.000199, \\
\rho_{\mathrm{clean}} &= +19.9\ \mathrm{pcm}, \\
x_{\mathrm{crit}} &\approx 0.193, \\
\Delta \rho_{\mathrm{rod}}(1.0) &= -164.8\ \mathrm{pcm}, \\
\rho_{\mathrm{scram}} &= -450\ \mathrm{pcm}\ \text{when latched.}
\end{align*}

The backend interpolates linearly between 11 precomputed insertion points
$x = 0.0, 0.1, \dots, 1.0$.

\subsection{Numerical method}

The point-kinetics equations are stiff, so the backend does not use explicit
Euler.  Instead \texttt{engine.py} uses a \textbf{semi-implicit elimination of
the precursors}.  For each step $\Delta t$,
\[
C_i^{k+1}
=
\frac{C_i^k + \Delta t\,(\beta_i/\Lambda)\,n^{k+1}}
{1 + \lambda_i \Delta t},
\]
and substituting that into the neutron balance yields a stable scalar update for
$n^{k+1}$.

In the implemented solver this is written in the form
\[
n^{k+1}
=
\frac{n^k + \Delta t\sum_i \dfrac{\lambda_i C_i^k}{1+\lambda_i\Delta t}}
{1 - \Delta t\,\dfrac{\rho-\beta}{\Lambda}
 - \Delta t\sum_i
   \dfrac{\lambda_i \Delta t\,\beta_i/\Lambda}{1+\lambda_i\Delta t}},
\]
which matches the code variables
\texttt{delayed\_source}, \texttt{delayed\_coupling}, \texttt{numerator}, and
\texttt{denominator} in \texttt{engine.py}.  The precursor update then uses the
newly computed $n^{k+1}$.

Numerically, this matters because the prompt part of the kinetics is very fast
relative to the slowest delayed groups.  A naive explicit step would require a
much smaller $\Delta t$ for stability, while the semi-implicit update remains
well behaved for the chosen browser-friendly step size.

\subsection{How the Overview transient is operationalized}

The point-kinetics equations are not advanced continuously in a background
thread.  Instead \texttt{service.py} advances them \emph{opportunistically}
whenever the frontend polls or sends a command:
\begin{enumerate}
  \item compute elapsed wall time since the last request,
  \item cap that wall-time gap to avoid large jumps after browser pauses,
  \item multiply by the configured time-scale factor,
  \item then march the engine forward in repeated $0.02$ s internal steps.
\end{enumerate}

This architecture is a numerical simplification as much as a software one: the
browser sees a continuous transient, but the backend actually performs a sequence
of bounded substeps and samples history only every $0.25$ s.

The current runtime tuning is:
\begin{center}
\begin{tabular}{lc}
\toprule
Parameter & Value \\
\midrule
Integrator step & $0.02$ s \\
History sampling & $0.25$ s \\
Time scale & $8\times$ wall clock \\
Frontend poll interval & $100$ ms \\
Automatic SCRAM threshold & $26$ MW$_{\mathrm{th}}$ \\
\bottomrule
\end{tabular}
\end{center}

\section{Core page: steady 2D diffusion solution}

The Core page is the repository's \textbf{spatial steady-state neutronics view}.
It reuses the same annular one-group model as the notebook and solves the 2D
eigenvalue problem for the \emph{current rod position} taken from the Overview
simulation.

\subsection{Discretization}

The implementation in \texttt{diffusion.py} and \texttt{core\_service.py} uses:
\begin{itemize}
  \item a cell-centered finite-difference mesh in $r$ and $z$,
  \item harmonic averaging of diffusion coefficients at material interfaces,
  \item sparse matrix assembly for leakage, absorption, and fission terms,
  \item and a power-iteration style eigenvalue solve for $k_{\mathrm{eff}}$ and
        the fundamental flux shape.
\end{itemize}

The production Core page solve runs at
\[
\Delta r = \Delta z = 3\ \mathrm{cm},
\]
and results are cached by rod insertion state so repeated requests do not repeat
the full eigenvalue solve.

\subsection{How the steady diffusion matrices are built}

The Python solver does not discretize the Laplacian in Cartesian form and then
``pretend'' the reactor is cylindrical.  It assembles the operator directly in
cylindrical coordinates.  On each cell-centered mesh point $(r_i,z_j)$ it builds:
\begin{itemize}
  \item radial couplings weighted by the cylindrical face areas $r_{i\pm 1/2}$,
  \item axial couplings weighted by $1/\Delta z^2$,
  \item and a diagonal absorption term from the local $\Sigma_a$.
\end{itemize}

Material interfaces are handled with the harmonic mean of adjacent diffusion
coefficients, which is the standard finite-volume-like choice for discontinuous
media because it preserves current continuity better than a simple arithmetic
average.

The assembled linear system can be written schematically as
\[
L\phi = \frac{1}{k}F\phi,
\]
where:
\begin{itemize}
  \item $L$ is a sparse loss matrix containing leakage and absorption,
  \item $F$ is a diagonal sparse matrix containing $\nu\Sigma_f$ in fuel cells
        and zero elsewhere.
\end{itemize}

\subsection{How the steady eigenvalue solve works numerically}

The function \texttt{solve\_2d()} uses a source-iteration / power-iteration
structure:
\begin{enumerate}
  \item start from a positive flux guess $\phi$ and an initial $k=1$,
  \item compute the fission source $F\phi$,
  \item solve the fixed loss system
        \[
        L\phi^{*} = F\phi
        \]
        with a sparse direct solver,
  \item clip tiny negative roundoff values to zero,
  \item normalize $\phi^{*}$ by its peak,
  \item update the eigenvalue with the ratio of new to old total fission source,
        \[
        k_{\mathrm{new}} = \frac{\sum(F\phi^{*})}{\sum(F\phi)},
        \]
  \item and repeat until $|k_{\mathrm{new}}-k| < \texttt{tol}$ after a few warmup
        iterations.
\end{enumerate}

Two numerical choices are especially important:
\begin{itemize}
  \item the loss matrix $L$ is factorized once per solve with
        \texttt{scipy.sparse.linalg.factorized()}, so each iteration is just a
        triangular solve instead of a fresh matrix inversion;
  \item the flux is normalized by peak value, not by integral power, because the
        steady eigenproblem only needs the shape and eigenvalue.
\end{itemize}

\subsection{How geometry and rod-worth calibration are solved}

The notebook and backend use the steady diffusion solver in two further numerical
loops:
\begin{enumerate}
  \item \textbf{critical-radius search:} bisection on $R_{\mathrm{fuel}}$ until
        the solved $k_{\mathrm{eff}}$ brackets unity to within tolerance;
  \item \textbf{rod-worth scan:} repeated solves at insertion fractions
        $x=0,0.1,\dots,1.0$, converting each $k$ to reactivity
        $\rho = (k-1)/k$ in pcm.
\end{enumerate}

The rod-worth scan may warm-start each solve with the previous flux shape.  That
is a purely numerical acceleration: neighboring rod positions have similar flux
shapes, so reusing the previous converged state reduces iteration count without
changing the underlying model.

\subsection{What the page returns}

The backend returns:
\begin{itemize}
  \item the peak-normalized 2D flux field $\phi(r,z)$,
  \item a radial profile through the midplane,
  \item an axial profile through the fuel-annulus midpoint,
  \item and metadata including $k_{\mathrm{eff}}$, mesh spacing, and iteration count.
\end{itemize}

This is the web app's direct counterpart to the notebook's Estimate 2 flux maps.

\section{Transient Diffusion page: direct time-dependent diffusion}

\subsection{Why this page exists}

The Overview page assumes a fixed spatial shape and evolves only an amplitude.
The Transient Diffusion page removes that simplification and evolves the
\textbf{spatial flux field itself}, together with the six delayed-neutron
precursor fields.

\subsection{Governing equations}

The transient diffusion solver in
\texttt{backend/reactor\_backend/transient\_diffusion.py} solves
\begin{align}
\frac{1}{v}\frac{\partial \phi}{\partial t}
&=
-L\phi
+ (1-\beta)\frac{F}{k_{\mathrm{ref}}}\phi
+ \sum_{i=1}^{6}\lambda_i C_i,
\\
\frac{\partial C_i}{\partial t}
&=
\beta_i \frac{\nu\Sigma_f}{k_{\mathrm{ref}}}\phi
- \lambda_i C_i,
\qquad i=1,\dots,6,
\end{align}
where:
\begin{itemize}
  \item $L$ is the spatial loss operator (leakage plus absorption),
  \item $F$ is the diagonal fission-production operator,
  \item $C_i(r,z,t)$ is the precursor field for group $i$,
  \item $v = 2.2\times 10^5\ \mathrm{cm/s}$ is the one-group thermal neutron speed,
  \item $k_{\mathrm{ref}}$ is the eigenvalue of the initialization state.
\end{itemize}

The division by $k_{\mathrm{ref}}$ is deliberate: it references the transient
fission source to the eigenvalue of the initialized state on the transient mesh.
That keeps the transient solve anchored to the same discrete criticality
reference and substantially reduces startup mismatch relative to using the raw
unscaled fission operator.

\subsection{Numerical method}

The transient page uses \textbf{fully implicit Euler}.  Eliminating the
precursors gives one sparse linear solve per time step:
\[
A\phi^{n+1} = b,
\]
with
\[
A = \frac{1}{v\Delta t}I + L
- b_{\mathrm{eff}}\,\mathrm{diag}\!\left(\frac{\nu\Sigma_f}{k_{\mathrm{ref}}}\right),
\]
and
\[
b_{\mathrm{eff}}
=
(1-\beta) + \sum_i \frac{\beta_i \lambda_i \Delta t}{1+\lambda_i \Delta t}.
\]

After $\phi^{n+1}$ is solved, each precursor field is updated by
\[
C_i^{n+1}
=
\frac{C_i^n + \Delta t\,\beta_i(\nu\Sigma_f/k_{\mathrm{ref}})\phi^{n+1}}
{1+\lambda_i \Delta t}.
\]

This scheme is unconditionally stable for the linear model and is much better
suited than explicit Euler to delayed-neutron reactor kinetics.

In implementation terms, the transient solver performs three numerically distinct
operations:
\begin{enumerate}
  \item \textbf{Initialization:} run a steady \texttt{solve\_2d()} eigenvalue
        solve on the transient mesh, flatten the 2D flux field, and initialize
        each precursor field with the steady relation
        \[
        C_{i,0} = \frac{\beta_i}{\lambda_i}\frac{\nu\Sigma_f}{k_{\mathrm{ref}}}\phi_0.
        \]
  \item \textbf{Matrix rebuild on rod motion:} if rod insertion changes, rebuild
        the sparse matrix $A$ and refactorize it.  In code this cache is keyed by
        the rounded rod fraction together with $k_{\mathrm{ref}}$.
  \item \textbf{Cheap repeated stepping between rod changes:} if rod insertion is
        unchanged, reuse the cached LU factorization so each new time step is
        just one sparse linear solve plus explicit vector updates for the six
        precursor groups.
\end{enumerate}

This split is why the transient page can remain interactive even though it is a
full spatial solver: the expensive part is the factorization, not the individual
post-factorization time steps.

\subsection{How the transient physics is manipulated into a solvable model}

The transient page is more detailed than point kinetics, but it is still not a
full reactor multiphysics model.  Several physics reductions are built into the
solver:
\begin{itemize}
  \item prompt and delayed neutrons are treated in one energy group,
  \item thermal-hydraulic feedback is absent, so all cross sections stay fixed
        while the transient runs,
  \item rod motion changes only the absorber map inside the central control-bank
        cylinder,
  \item the internally stored flux is an \emph{absolute amplitude field}, while
        UI heatmaps are separately peak-normalized display copies,
  \item and power is derived from a cylindrical fission integral rather than from
        a standalone thermal-hydraulic power balance.
\end{itemize}

That last point is numerically important: the solver never renormalizes its
internal state after each step.  If it did, the transient could not display real
growth or decay in $P/P_0$.

\subsection{Runtime settings}

To keep the direct transient responsive in a browser-driven local app, the
production solver uses:

\begin{center}
\begin{tabular}{lc}
\toprule
Parameter & Value \\
\midrule
Transient mesh & $\Delta r = \Delta z = 5$ cm \\
Time step & $\Delta t = 1$ s \\
Display heatmap & $40 \times 60$ points \\
History limit & 200 points \\
\bottomrule
\end{tabular}
\end{center}

The transient mesh is intentionally coarser than the Core page mesh so the LU
factorization can be rebuilt quickly whenever rod insertion changes.

\subsection{Service-level time stepping}

The transient page is advanced by \texttt{transient\_service.py}, which keeps one
shared solver instance and one shared state vector.  On every
\texttt{GET /api/transient-diffusion/state} request it:
\begin{enumerate}
  \item measures elapsed wall time,
  \item converts that to simulated time at a fixed scale of $1\times$,
  \item computes how many whole 1 s transient steps fit in that interval,
  \item caps the advance to at most four steps per request,
  \item and appends one history sample per accepted time step.
\end{enumerate}

This cap is not part of the physics model; it is a numerical and UX safeguard so
one delayed frontend poll cannot trigger an unexpectedly long blocking solve.

\subsection{Power normalization}

The transient diffusion page no longer assumes that power is proportional to one
global amplitude alone.  Instead it evaluates the cylindrical fission-rate
integral
\[
P \propto \sum_j \nu\Sigma_{f,j}\,\phi_j\,V_j,
\qquad
V_j \propto r_j \Delta r \Delta z,
\]
and reports
\[
\frac{P}{P_0}
=
\frac{\sum_j \nu\Sigma_{f,j}\phi_j V_j}
{\sum_j \nu\Sigma_{f,j}\phi_{j,0} V_j}.
\]

The factor $2\pi$ cancels in the ratio, so the implementation only needs the
$r_j \Delta r \Delta z$ weight per cell.

\section{How the three app pages relate}

\begin{center}
\begin{tabular}{p{2.6cm}p{3.8cm}p{6.1cm}}
\toprule
Layer & Model & Purpose \\
\midrule
Notebook estimate 1 & Homogeneous buckling & First-pass analytical check of scale and leakage trends \\
Notebook estimate 2 / Core page & Steady 2D one-group diffusion & Compute $k_{\mathrm{eff}}$, flux maps, and rod-worth calibration \\
Overview page & 0D six-group point kinetics & Fast transient driven by the calibrated reactivity curve \\
Transient Diffusion page & 2D diffusion + 6 precursor fields & Spatial transient with evolving flux shape and power \\
\bottomrule
\end{tabular}
\end{center}

The conceptual workflow is therefore:
\begin{enumerate}
  \item derive or justify effective constants in the notebook,
  \item solve the steady annular diffusion problem and calibrate rod worth,
  \item freeze those results into backend constants,
  \item use the fast point model for dashboard-style operation,
  \item and use the diffusion transient when spatial dynamics matter.
\end{enumerate}

\section{Important assumptions and limitations}

The current model remains intentionally simplified:
\begin{itemize}
  \item one energy group only,
  \item no thermal-hydraulic feedback,
  \item no fuel temperature, moderator temperature, void, xenon, or burnup effects,
  \item homogenized fuel and rod regions,
  \item heavy-water moderation is represented through effective diffusion and
        absorption constants rather than coupled moderator state equations.
\end{itemize}

So the repository should be read as an \textbf{educational reactor simulator with
increasing levels of neutronics fidelity}, not as a licensing-grade core code.

\section{Bottom line}

\texttt{reactorModel.ipynb} is the repository's physics workbench: it derives the
effective one-group annular-core picture, performs the 2D criticality and
rod-worth studies, and defines the constants later frozen into the backend.

The web app then uses that physics in two ways:
\begin{itemize}
  \item as a reduced \textbf{point-kinetics transient} for fast control-response
        exploration on the Overview page,
  \item and as a direct \textbf{diffusion-based spatial solve}, first in steady
        form on the Core page and then in time-dependent form on the Transient
        Diffusion page.
\end{itemize}

That is the current modeling hierarchy of the project.

\end{document}
